\DocumentMetadata{
  pdfversion=2.0,
  pdfstandard=ua-2,
  lang=en-US,
  tagging=on,
  tagging-setup={math/setup=mathml-SE,tabsorder=structure}
}
\documentclass[11pt,letterpaper]{article}
\usepackage[margin=0.68in,includefoot]{geometry}
\usepackage{newcomputermodern}
\usepackage{amsmath,amscd,mathtools}
\usepackage{tikz,adjustbox}
\providecommand{\square}{\mdlgwhtsquare}
\usepackage[hidelinks]{hyperref}
\hypersetup{
  pdftitle={Numerical Analysis First-Year Exam, August 2020},
  pdfauthor={University of Florida Department of Mathematics},
  pdfsubject={Graduate mathematics examination},
  pdfdisplaydoctitle=true,
  bookmarksopen=true
}
\setcounter{secnumdepth}{0}
\setlength{\parindent}{0pt}
\setlength{\parskip}{0.28em}
\setlength{\emergencystretch}{3em}
\setlength{\leftmargini}{2.15em}
\setlength{\leftmarginii}{2.65em}
\setlength{\leftmarginiii}{3em}
\renewcommand{\labelenumi}{\textbf{\theenumi.}}
\pagestyle{empty}
\raggedbottom
\allowdisplaybreaks
\newenvironment{examproblems}[1]{%
  \begin{enumerate}%
  \setcounter{enumi}{#1}%
  \setlength{\topsep}{0.35em}%
  \setlength{\partopsep}{0pt}%
  \setlength{\itemsep}{0.48em}%
  \setlength{\parsep}{0.18em}%
}{\end{enumerate}}
\newenvironment{examsubparts}{%
  \begin{enumerate}%
  \setlength{\topsep}{0.18em}%
  \setlength{\partopsep}{0pt}%
  \setlength{\itemsep}{0.16em}%
  \setlength{\parsep}{0.1em}%
}{\end{enumerate}}
\newenvironment{examinstructions}{%
  \par\begingroup\itshape
}{%
  \par\endgroup\vspace{0.18em}
}
\makeatletter
\renewcommand\section{\@startsection{section}{1}{0pt}%
  {-1.25ex plus -.25ex minus -.1ex}%
  {0.55ex plus .1ex}%
  {\normalfont\large\bfseries}}
\renewcommand{\@maketitle}{%
  \newpage\null\vskip -1em%
  {\footnotesize\bfseries
    \href{https://www.ufl.edu/}{University of Florida}, \href{https://math.ufl.edu/}{Department of Mathematics}\par}%
  \vskip -0.5em%
  \tagstructbegin{tag=Title}%
    {\LARGE\bfseries\href{https://gma.math.ufl.edu/exams/numerical-analysis-first-year/}{Numerical Analysis First-Year Exam}, August 2020\par}%
  \tagstructend%
  \vskip 0.25em%
  \tagmcbegin{artifact}%
    \hrule height 1pt%
  \tagmcend%
  \vskip 1em%
}
\makeatother
\title{Numerical Analysis First-Year Exam, August 2020}
\author{}
\date{}
\begin{document}
\maketitle
\thispagestyle{empty}
\begin{examinstructions}
Do 4 (four) problems.
\end{examinstructions}
\begin{examproblems}{0}
\item
Let \(\{\phi_k\}_{k=0}^{n+1}\) be a set of orthogonal polynomials on \([a,b]\), with respect to the inner product \((f,g)=\int_a^b f(x)g(x)w(x)\,dx\), indexed so that \(\phi_k\) is of degree~\(k\). Prove that \(\phi_k\) has~\(k\) distinct roots \(\{x_j^{(k)}\}_{j=1}^k\), with \(x_j^{(k)}\in[a,b]\), \(j=1,\ldots,k\).
\item
Consider the interval \([a,b]\) with the partition \(a=x_1<x_2<\cdots<x_n<x_{n+1}=b\). Suppose \(s(x)\) is the natural cubic spline that interpolates the data \(\{(x_i,y_i)\}_{i=1}^{n+1}\), and that \(g\in C^2[a,b]\) interpolates the same data. Show that 
\[
\int_a^b (s''(x))^2\,dx\leq\int_a^b (g''(x))^2\,dx.
\]
\item
This problem has two parts.
\begin{examsubparts}
\item[\textnormal{(a)}]
Consider the inner product on \(C(0,2)\) given by \((f,g)=\int_0^2 f(t)g(t)\,dt\). Find three orthonormal polynomials \(\phi_0,\phi_1,\phi_2\) on \((0,2)\) with respect to the given inner product such that the degree of \(\phi_n\) is equal to~\(n\), \(n=0,1,2\).
\item[\textnormal{(b)}]
Find the nodes \(t_1\) and \(t_2\) and weights \(w_1\) and \(w_2\) which yield the weighted Gaussian Quadrature formula 
\[
\int_0^2 f(t)\,dt\approx w_1f(t_1)+w_2f(t_2)
\]
 with degree of exactness \(m=3\). You should find the nodes exactly, and may leave the weights \(w_1,w_2\) in integral form.
\end{examsubparts}
\item
Prove for any \(f\in C[a,b]\) and integer \(n\geq 0\), that the best uniform approximation of~\(f\) in \(P_n\) is unique. You may assume the existence of at least one best uniform approximation of~\(f\).
\item
Suppose \(f\in C^{n+1}[a,b]\), and let \(p\in P_n\) be a polynomial that interpolates the data \(\{(x_i,f(x_i))\}_{i=0}^n\), where \(x_0,\ldots,x_n\) are distinct points in \([a,b]\). Consider an arbitrary fixed \(x\in[a,b]\), and derive an exact expression for the error \(f(x)-p(x)\).
\end{examproblems}
\end{document}
